% Part: first-order-logic
% Chapter: proof-systems
% Section: introduction

\documentclass[../../../include/open-logic-section]{subfiles}

\begin{document}

\iftag{FOL}
      {\olfileid{fol}{prf}{int}}
      {\olfileid{pl}{prf}{int}}

\olsection{ভূমিকা}

যুক্তিবিদ্যায় সাধারণত একটি অর্থতত্ত্ব ও !!a{derivation}
পদ্ধতি—দুটিই থাকে। অর্থতত্ত্ব সত্যতা, পরিতৃপ্তিযোগ্যতা, বৈধতা ও
অর্থগত অনুসিদ্ধান্তের মতো ধারণা নিয়ে আলোচনা করে। !!{derivation}
পদ্ধতির উদ্দেশ্য হলো অনুসিদ্ধান্ত ও বৈধতা প্রতিষ্ঠার জন্য সম্পূর্ণ
সংকেতবিন্যাসগত একটি পদ্ধতি দেওয়া। এগুলি সম্পূর্ণ সংকেতবিন্যাসগত এই অর্থে
যে, এমন পদ্ধতিতে কোনো !!{derivation} একটি সসীম সংকেতবিন্যাসগত বস্তু;
সাধারণত !!{sentence}s বা !!{formula}s-এর একটি অনুক্রম (বা অন্য কোনো
সসীম বিন্যাস)। ভালো !!{derivation} পদ্ধতির বৈশিষ্ট্য হলো, !!{sentence}s
বা !!{formula}s-এর যেকোনো প্রদত্ত অনুক্রম বা বিন্যাস “সঠিক” কি না,
তা যান্ত্রিকভাবে যাচাই করা যায়।

প্রথম-ক্রমের যুক্তিবিদ্যার সবচেয়ে সরল (এবং ঐতিহাসিকভাবে প্রথম)
!!{derivation} পদ্ধতিগুলি ছিল \emph{স্বতঃসিদ্ধমূলক}। এমন পদ্ধতিতে
!!{formula}s-এর কোনো অনুক্রমকে !!a{derivation} বলা হয়, যদি তার
প্রতিটি !!{formula} হয় কোনো নির্দিষ্ট “স্বতঃসিদ্ধ”-সমষ্টির অন্তর্গত,
নয়তো নির্দিষ্টসংখ্যক “অনুমান-বিধি”-র কোনো একটির সাহায্যে অনুক্রমে তার
আগে থাকা !!{formula}s থেকে অনুসৃত হয়। কোনো !!a{formula} স্বতঃসিদ্ধ কি
না এবং কোনো অনুমান-বিধি দিয়ে অন্য !!{formula}s থেকে সেটি সঠিকভাবে
অনুসৃত হয় কি না, তা যান্ত্রিকভাবে যাচাই করা যায়। স্বতঃসিদ্ধমূলক
!!{derivation} পদ্ধতি বর্ণনা করা সহজ এবং অধিতাত্ত্বিকভাবে বিচার করাও
সহজ; কিন্তু তাতে !!{derivation}s পড়া ও বোঝা কঠিন, তৈরি করাও কঠিন।

অন্য !!{derivation} পদ্ধতি তৈরি হয়েছে, যাতে !!{derivation}s নির্মাণ
করা অথবা নির্মাণ-শেষের !!{derivation}s বোঝা সহজ হয়। উদাহরণ হলো স্বাভাবিক
নিষ্পাদন, সত্য-বৃক্ষ—যা ট্যাবলো প্রমাণ নামেও পরিচিত—এবং সিকোয়েন্ট কলন।
কিছু !!{derivation} পদ্ধতি বিশেষভাবে যান্ত্রিক প্রয়োগের কথা ভেবে
পরিকল্পিত; যেমন, রেজোলিউশন পদ্ধতি সফটওয়্যারে রূপায়ণ করা সহজ (কিন্তু তার
!!{derivation}s বোঝা কার্যত অসম্ভব)। এই অন্য !!{derivation} পদ্ধতিগুলির অধিকাংশই
!!{derivation}s-কে অনুক্রমের বদলে !!{formula}s-এর বৃক্ষ হিসেবে প্রকাশ
করে। ফলে !!a{derivation}-এর কোন অংশ অন্য কোন অংশের উপর নির্ভর করে,
তা দেখা সহজ হয়।

তাই প্রথম-ক্রমের যুক্তিবিদ্যার মতো কোনো যুক্তিবিদ্যার ক্ষেত্রে বিভিন্ন
!!{derivation} পদ্ধতি, !!a{sentence}-এর \emph{উপপাদ্য} হওয়া কী এবং
অন্য কয়েকটি থেকে কোনো !!a{sentence} !!{derivable} হওয়ার অর্থ কী—তার
আলাদা আলাদা স্পষ্টীকরণ দেবে। যে ভাবেই তা করা হোক (স্বতঃসিদ্ধমূলক
!!{derivation}s, স্বাভাবিক নিষ্পাদন, সিকোয়েন্ট !!{derivation}s,
সত্য-বৃক্ষ, রেজোলিউশন খণ্ডনের মাধ্যমে), আমরা চাই এই সম্পর্কগুলি বৈধতা
ও অর্থগত অনুসিদ্ধান্তের অর্থগত ধারণার সঙ্গে মিলে যাক। “$!A$~একটি
উপপাদ্য” বোঝাতে $\Proves !A$ এবং “$\Gamma$ থেকে $!A$~!!{derivable}”
বোঝাতে “$\Gamma \Proves !A$” লিখি। $\Proves$ যেভাবেই সংজ্ঞায়িত হোক,
আমরা চাই সেটি $\Entails$-এর সঙ্গে মিলে যাক; অর্থাৎ:
\begin{enumerate}
\item $\Proves !A$ যদি এবং কেবল যদি $\Entails !A$
\item $\Gamma \Proves !A$ যদি এবং কেবল যদি $\Gamma \Entails !A$
\end{enumerate}
উপরের “কেবল যদি” অভিমুখকে
\emph{বিশুদ্ধতা} বলা হয়। কোনো !!^a{derivation} পদ্ধতি বিশুদ্ধ, যদি
!!{derivability} অর্থগত অনুসিদ্ধান্তের (বা বৈধতার) নিশ্চয়তা দেয়।
প্রতিটি গ্রহণযোগ্য !!{derivation} পদ্ধতিই বিশুদ্ধ হতে হবে; অবিশুদ্ধ
!!{derivation} পদ্ধতি একেবারেই উপযোগী নয়। কারণ !!a{derivation}-এর
সমগ্র উদ্দেশ্যই হলো বৈধতা বা অর্থগত অনুসিদ্ধান্তের একটি সংকেতবিন্যাসগত
নিশ্চয়তা দেওয়া। আমাদের উপস্থাপিত !!{derivation} পদ্ধতিগুলির বিশুদ্ধতা
আমরা প্রমাণ করব।

বিপরীত “যদি” অভিমুখটিও গুরুত্বপূর্ণ: তাকে
\emph{পূর্ণতা} বলা হয়। কোনো পূর্ণ !!{derivation} পদ্ধতি এতটাই শক্তিশালী
যে $!A$ বৈধ হলেই সেটি $!A$-কে উপপাদ্য হিসেবে দেখাতে পারে, এবং
$\Gamma \Entails !A$ হলেই $\Gamma \Proves !A$ দেখাতে পারে।
পূর্ণতা প্রতিষ্ঠা করা কঠিনতর, এবং কিছু যুক্তিবিদ্যার কোনো পূর্ণ
!!{derivation} পদ্ধতি নেই। প্রথম-ক্রমের যুক্তিবিদ্যার আছে। কুর্ট গ্যোডেল
তাঁর ১৯২৯ সালের অভিসন্দর্ভে প্রথম-ক্রমের যুক্তিবিদ্যার একটি
!!a{derivation} পদ্ধতির পূর্ণতা প্রথম প্রমাণ করেন।

!!{derivation} পদ্ধতির সঙ্গে যুক্ত আরেকটি ধারণা হলো
\emph{সঙ্গতি}। কোনো !!{sentence}s-এর সেট থেকে যা কিছুই হোক
!!{derive}d করা গেলে সেটটিকে অসঙ্গত বলা হয়; অন্যথায় সেটি সঙ্গত।
অসঙ্গতি হলো অপরিতৃপ্তিযোগ্যতার সংকেতবিন্যাসগত প্রতিরূপ: অপরিতৃপ্তিযোগ্য
সেটের মতো অসঙ্গত !!{sentence}s-এর সেটও ভালো তত্ত্ব হয় না; তাদের মৌলিক
ত্রুটি থাকে। সঙ্গত !!{sentence}s-এর সেট সত্য বা উপযোগী না-ও হতে পারে,
কিন্তু অন্তত যৌক্তিক উপযোগিতার এই ন্যূনতম সীমাটি পার হয়। বিভিন্ন
!!{derivation} পদ্ধতিতে !!{sentence}s-এর সেটের সঙ্গতির নির্দিষ্ট সংজ্ঞা
আলাদা হতে পারে। তবে~$\Proves$-এর মতো সঙ্গতিও তার অর্থগত প্রতিরূপ
পরিতৃপ্তিযোগ্যতার সঙ্গে মিলে যাক—আমরা সেটিই চাই। আমরা চাই, সব সময়
$\Gamma$ সঙ্গত যদি এবং কেবল যদি সেটি পরিতৃপ্তিযোগ্য। এখানে “কেবল যদি”
অভিমুখটি পূর্ণতার সমতুল্য (সঙ্গতি পরিতৃপ্তিযোগ্যতার নিশ্চয়তা দেয়), এবং
“যদি” অভিমুখটি বিশুদ্ধতার সমতুল্য (পরিতৃপ্তিযোগ্যতা সঙ্গতির নিশ্চয়তা
দেয়)। বস্তুত, ধ্রুপদি প্রথম-ক্রমের যুক্তিবিদ্যায় বিশুদ্ধতা ও পূর্ণতার
দুই রূপই পরস্পর সমতুল্য।

\end{document}
